\section{Routing}
Problema: Il nodo x vuole comunicare un'informazione con l'entità y. Come fa x a raggiungere y?\\
Broadcast routing: x manda in broadcast l'informazione a tutta la rete. Soluzione più semplice, funziona sempre, ma alto numero di
messaggi e problema di confidenzialità.\\
\textbf{Routing}: processo che consiste nel determinare un percorso tra sorgente x e destinazione y.\\
\textbf{Router}: nodo dedicato a determinare quale percorso deve seguire un messaggio per raggiungere la destinazione.

Ogni entità ha un \textbf{identificativo univoco} (global names).
Ogni link ha un'orientamento locale. Dei costi possono essere associati ai link.
Ogni entità ha una routing function $f()$. Per ogni destinazione $y$, $f(y)$ è il link sul quale il messaggio deve essere mandato per raggiungere $y$.

La decisione di routing è pressa sulla base di un'informazione presente nel messaggio (\textbf{indirizzo del destinatario}) e un'informazione locale (\textbf{routing table}).

Stiamo attenti a: routing memory, ovvero i bits necessari per memorizzare la tabella, e routing time, ovvero il tempo necessario per trovare il next hop.

Esempio: tabella con ogni riga contenente il nodo dst, il next hop ed, eventualmente il costo.
% TODO immagine.

Costo Temporale: $O(n \log n)$ con binary search siccome è ordinata.\\
Costo Spaziale: $O(n \log \text{max}_\text{id})$; n righe, ognuna con ID del next hop.

I cammini minimi rispettano il \textbf{principio di subottimalità}. Per questo possiamo ricordarci solo il next hop nella tabella di routing.

Soluzione semplice. Se ogni nodo conoscesse a priori tutto il grafo, potrebbe costruirsi l'albero dei cammini minimi a partire da quello e ricavare i next hop.

\subsection{Gossip}
Idea: ogni nodo deve arrivare ad avere una vista completa del grafo. \\
Esecuzione: ogni nodo comunica a tutti gli altri quali sono i suoi vicini e quanto costano gli archi verso di loro.
Broadcast tutti a tutti:
\begin{itemize}
  \item contruiamo uno spanning tree qualunque per limitare il numero dei messaggi. $O(2m)$ con SHOUT+.
  \item ogni entità riceve le informazioni dal suo vicino. Ogni coppia di nodi si scambia info, quindi $O(2m)$.
  \item ogni entità propaga l'info sullo ST. Ogni entità $x$ sui suoi vicini nel grafo ($N(x)$) a tutti i suoi vicini in ST: $\sum_x (n-1) \text{deg}(x) = (n-1) \sum_x \text{deg}(x) = (n-1) (2m)$
  \item ogni entità calcola localmente lo stortest path. Questa operazione non richiede messaggi.
\end{itemize}

Dai costi sopra, il totale è $\text{Messages}(\text{Gossip}(G)) = O(mn)$.

Problema: non è detto che tutte le entità abbiano spazio sufficiente per memorizzare l'intero grafo\dots cerchiamo una soluzione più efficiente.

\subsection{Bellman Ford Distribuito}
Idea: ogni entità conosce i suoi vicini e iterativamente aggiorna il suo \textbf{distance vector} man mano che i vicini gli propagano informazioni
su nodi sempre più lontani.

Ogni riga della tabella di routing può contenere il next hop con un costo associato (in un dato momento prima della convergenza,
potenzialmente subottimale), oppure un unknown next hop con un costo infinito, perchè irraggiungibile.

Questo protocollo è una implementazione distribuita di $Bellman-Ford$. Richiede al massimo $n-1$ iterazioni (lunghezza massima di un cammino
minimo su grafo) prima di arrivare a convergenza. È usato in protocolli di routing come RIP. Durante le iterazioni, ogni entry della routing
table viene sostituita se, alla luce del distance vector ricevuto, è possibile abbassare il costo per raggiungere uno dei nodi della tabella.

Message Complexity. Ad ogni iterazione ogni entità manda il proprio distance vector ai propri vicini.
Ogni messaggio contiene $n$ righe con 2 colonne, con destinazione e costo. $O(n)$ propagato a $|N(x)|$ nodi.

Messaggi per iterazione: $n \sum_x |N(x)|$ = $n \cdot 2m$\\
Messaggi totali: $\text{Messages}(\text{BellmanFordDistributed}(G)) = (n-1) \cdot n \cdot 2m \in O(m n^2)$

Idea: possiamo evitare di costruire una vista globale per ogni entità? sì, possiamo adottare un approccio distribuito.

Prima però ragioniamo sull'aggiornamento delle tabelle di routing in caso di aggiornamento dei nodi. Come gestire i cambiamenti della topologia della rete?

Map update: un cambiamento nella rete rende le mappe locali inaccurate.\\
Distance vector update: un cambiamento nella rete potrebbe portare all'aggiornamento di uno o tutti (nel caso limite in cui cambia un link
da cui passa tutto il traffico della rete) i distance vector.

Quando una entità identifica un cambiamento locale possiamo:
\begin{itemize}
 \item ricalcolare tutte le routing tables (dispendioso)
 \item aggiornare le tabelle di routing correnti. Ogni entità inizia una nuova iterazione (\textit{update} message) del protocollo usando il nuovo dato, finchè la tabelle convergono
\end{itemize}

Quando un'entità riceve un \textit{update}: aggiorna la sua routing table, propaga l'aggiornamento.
La rete potrebbe non essere ancora arrivata a convergenza prima del prossimo aggiornamento della rete.

Mantenere le routing table costantemente aggiornate con il shortest path aggiunge overhead alla rete. Nel mondo reale si preferisce
\textbf{rilassare} il vincolo di shortest path, lasciando solo quello di \textbf{delivery}.

\subsection{Min-Hop Routing}
Assuzione: tutti i link hanno lo stesso costo, non memorizziamo i costi, li assumiamo \textbf{unitari}.

Il Breadth First Spanning Tree (BFST) è lo Shortest Path Spanning Tree (SPST) in questo caso. BFST può essere calcolato con SHOUT nel caso sincrono.
Nel caso asincorno non abbiamo questa garanzia: potremmo avere link molto più veloci di altri, nel caso estremo una catena di link che
portano prima in profondità che in larghezza.

Come facciamo su sistemi \underline{asincroni}?

\obsbox{
  Un nodo al \textbf{livello} $i$ avrà vicini solamente alla distanza $i$ (fratelli), $i-1$ (almeno uno, il padre), e $i+1$ (figli).
  Nei sistemi asincroni dobbiamo inglobare i nodi per livelli, dove ogni livello è determinato dalla distanza dalla radice.
  È impossibile che un nodo al livello $i$ abbia vicini al livello $i-2$, altrimenti il nodo in questione sarebbe al livello $i-1$.
  Per determinare i nodi alla distanza $i+1$ sono necessari solamente i vicini dei nodi alla distanza $i$.
}

Possiamo iniziare a costrurie il prossimo livello solamente quando quello corrente è stato completato per intero.
I vicini dei nodi del livello corrente potrebbero essere già nell'albero se appartengono al livello precedente o sono nodi fratelli.

L'orchestratore e inziatore di questo processo sarà la \textbf{radice}.
Stiamo relalizzando la versione distribuita del \textbf{BFST}. Costruiamo iterativamente lo ST.
Per comunicare tra i nodi sfruttiamo l'albero già costruito.

Alla generatica iterazione $i$ il \textbf{partial tree} contiene i nodi fino al livello $i-1$; vogliamo inglobare i nodi del livello $i$.
Abbiamo 3 stadi:
\begin{enumerate}
  \item la radice propaga un messaggio di ``inizio dell'iterazione i'' al partial tree.
  \item i nodi del livello $i-1$ (conoscono a quale livello si trovano, dato che corrisponde all'iterazione a cui sono stati inclusi nel partial tree)
  avviano la procedura \textit{Exploring}; i nodi a livelli $< i-1$ si limitano a propagare il messaggio ai figli nel partial tree.\\
  Finito \textit{Exploring} i nodi del livello $i-1$ avvisano il padre con \textit{iteration i ended}.
  \item i nodi dei livelli $< i-1$ attendono la terminazione di tutti i figli e propagano \textit{iteration i ended} al padre.
  Se tutti i figli della radice finiscono, si inizia con l'iterazione successiva.
\end{enumerate}

\textit{Exploring}:
\begin{itemize}
  \item ogni entità x alla distanza $i-1$:
  \begin{enumerate}
    \item invia ``explore i'' ai ai vicini diretti, escluso il genitore del nodo nell'albero
    \item aspetta un ack da ogni nodo vicino contattato. Se riceve:
    \begin{itemize}
      \item ack-OK il vicino diventa un figlio
      \item ack-KO il vicino è già nell'albero (lv $i-2$, lv $i-1$, o lv $i$ ma già contattato ad un altro nodo)
      \item explore-i il vicino è già nell'albero, al livello $i-1$. Vedrà anch'esso il nostro explore-i.
    \end{itemize}
  \end{enumerate}
  \item entità y non appartenente all'albero, alla ricezione di ``explore i'':
  \begin{itemize}
    \item la prima volta invia ack-OK e segna il nodo da cui ha ricevuto explore i come parent. y si trova a distanza $i$.
    \item dalla volta dopo in poi invia ack-KO
  \end{itemize}
\end{itemize}
    
Come terminare la procedura?\\
La radice si deve accorgere che abbiamo incluso tutti i nodi e inviare una notifica di terminazione.
Un'idea funzionante ma dispendiosa è fare al più $N-1$ iterazioni (caso peggiore, nodi in catena), conoscendo $N$ numero di nodi.

Un'alternative più ottimizzata consiste nello sfruttare la fase di terminazione.
La fase di invio dei messaggi di terminazione è un \textbf{convergecast}. Il \textbf{convergecast} può essere esteso aggiungendo l'informazione
di assenza di ulteriori nodi su cui espandere l'albero.
Un nodo del livello corrente sa di essere una foglia se riceve solo ack-KO dai vicini, quindi propaga al genitore un valore speciale.
Un generico nodo, propaga il valore speciale verso la radice solamente se ha ricevuto solo valori speciali dai figli, a.k.a i suo i figli sono
tutte foglie.

Message complexity.\\
Il raggio di una sorgente è la distanza massima delle foglie dell'albero BFST dalla radice, ovvero $r(s) = \text{altezza di BFST radicato in s}$.\\
Alla generica iterazione $i$:
\begin{enumerate}
 \item \textbf{Broadcast} al partial tree fino al livello $i-1$:
 risulta che $n_i - 1$ (conto il numero di link dell'albero parziale). Con $n_k$ nodi dell'albero parziale all'inizio dell'iterazione $k$.
 \item \textbf{Explore}. Ogni nodo manda esattamente 1 explore; non è possibile lo mandi più di una volta perchè vorrebbe dire che è stato considerato due volte per
 lo stesso livello. Su ogni link su cui è stato mandato un explore tornerà indietro un ack-OK, un ack-KO o un altro explore-i.
 \item \textbf{Convergecast} per propagare ``iterazione i finita''. Stesso numero di (1) dato che i messaggi devono tornare indietro per gli stessi link.
\end{enumerate}

Sommando (1) e (3):
$$
\sum_{i=1}^{r(s)} 2 (n_i - 2) \leq \sum_{i=1}^{r(s)} 2 (n - 2) = 2 (n - 2) \sum_{i=1}^{r(s)} 1 = 2 (n - 2) r(s) \leq 2 (n-2) D(G)
$$

Per (2): ogni link vede esattamente 2 messaggi, totalizzando $2m$ messaggi.

Totale:
$$
\text{Messages}(\text{MinHopRouting}(G_\text{async})) = 2(n-1)D(G) + 2m \leq 2(n-1)(n-1)+2m \in O(n^2)
$$

Si può migliorare in termini di messaggi?\\
Sì, posso evitare di mandare il messaggio di explore ai nodi del livello prima;
in questo modo il NO-ACK viene usato per segnalare che il nodo è già stato collegato ad il nodo padre.

Miglioramento? Ogni arco viene coinvolto una sola volta nella procedura explore-risposta.
Solo due messaggi per arco nella fase di explore: 2m.

Time complexity.\\
Alla generica iterazione $i$:
\begin{enumerate}
  \item \textbf{broadcast} nell'albero parziale, fino alla distanza $i-1$: catena di $i-1$ messaggi
  \item \textbf{explore}: catena lunga 2 messaggi
  \item \textbf{convergecast}: di nuovo una catena di $i-1$ messaggi
\end{enumerate}

Sommando tra due iterazioni:
$$
\text{Time}(\text{MinHopRouting}(G_\text{async})) = \sum_{i=1}^{r(s)} 2(i-1) + 2 = \sum_{i=1}^{r(s)} 2i =
2 \sum_{i=1}^{r(s)} i = \frac{r(s) \cdot [r(s)+1]}{2} \in O(r(s)^2) \in O(n^2)
$$

dato che il diametro peggiore è n, ovvero una catena.

\subsection{Shortest Path Spanning Tree}
Reintroduciamo ora i pesi degli archi.
Restrizione aggiuntiva: peso degli archi \textbf{positivo}.

Realizziamo la versione distribuita di \underline{Dijkstra}. Ricorda che è un algoritmo greedy, che fa uso di una coda di priorità.
La soluzione geedy viene tradotta con una soluzione distribuita che lavora per diverse iterazioni; come prima, la sorgente agisce da coordinatore.
Alla generica iterazione $i$ abbiamo una parte di albero dei cammini minimi formato e una parte dell'albero originale ancora da esplorare.

% TODO diagram.

Definiamo \textit{outgoing edge} (outgoing dalla propsettiva dello SPST) come un arco che collega il SPST parziale verso un nodo non ancora considerato;
si trova quindi alla frontiera dello SPST. Ad ogni iterazione, aggiungiamo un outgoing edge e il nodo su cui incide allo SPST.

Inizializzazione:
\begin{itemize}
 \item tutti gli archi sono impostati ad outgoing
 \item tutte le entità conoscono il costo degli archi incidenti
 \item solo la sorgente si trova nell'albero, tutte le altre entità sono fuori
 \item la sorgente notifca di essere all'interno dell'albero e attende un ack
 \item i vicini della sorgente impostano l'arco che li connette ad essa come non outgoing (dalla loro propsettiva; per la sorgente resta outgoing).
\end{itemize}

Posta $\Delta(x)$ minima distanza dalla sorgente $s$ al generico nodo $x$. Posto $c(x, y)$ costo dell'arco $(x, y)$.

Generic iteration:
\begin{itemize}
 \item La radice propaga \textit{start iteration} nello SPST parziale. Costo $O(n)$ messaggi.
 \item Una entità generica $x$ nel SPST parziale, alla ricezione di \textit{start iteration}:
 \begin{itemize}
    \item calcola la somma tra il costo per arrivare a $x$ e il costo di aggiunta di ogni nodo: $\Delta(x) + c(x, y) \quad \forall (x, y) \in \text{outgoing edges}$
    \item seleziona l'arco $(x, y_x)$ che minimizza tale costo (trovando il minimo locale)
 \end{itemize}
 \item L'infromazione di costo e arco aggiunto fluisce verso la radice mediante convergecast. Ogni nodo contribuisce al calcolo di $\hat x = \arg \min_x \Delta(x) + c(x, y_x)$.
 A.k.a., propaga verso l'alto solo il messaggio con costo minore tra quelli ricevuti dai figli. Costo $O(n)$ siccome il messaggio deve essere propagato sull'intero
 SPST ($n-1$ nodi nel caso peggiore).
 \item La radice aggiunge $\hat x$ e $(\hat x, y_x)$ allo SPST.
 \item La radice usando lo SPST notifica $y_{\hat x}$ della sua inclusione. $O(n)$ messaggi. Durante questa fase, $\hat x$ invia a $y_{\hat x}$ il valore
 $\Delta(y_{\hat x})$, ovvero il costo per raggiungere $y_{\hat x}$ dalla sorgente; l'arco $\hat x, y_{\hat x}$ viene marcato come appartenente a SPST.
 \item $y_{\hat x}$ notifica i vicini della sua inclusione e attente un ack. $O(\text{deg}(y_{\hat x}))$ messaggi.
 \item I vicini impostano l'arco verso $y_{\hat x}$ come non-outgoing e inviano un ack. $O(\text{deg}(y_{\hat x}))$ messaggi.
 \item Non appena $y_{\hat x}$ riceve tutti gli ack invia \textit{end iteration} alla radice. $O(n)$ messaggi.
\end{itemize}

Terminazione: per costruzione, la radice invia una notifica. Come?\\
La radice conosce $n$ e sa che, effettuando$n-1$ iterazioni tutti i nodi verranno considarati; vengono potenzialmente fatte più iterazioni del necessario.
Un'alternativa più efficiente consiste nel far inviare un valore speciale (e.g., negativo) ai nodi senza archi outgoing, durante il convergecast (valori chiaramente
esclusi dal calcolo del nodo di costo minimo); quando riceve solo valori speciali, sa che sono stati esauriti gli archi inesplorati.

Ricorda: ogni iterazione un solo nodo è aggiunto al SPST quindi, nel caso peggiore, serviranno $n$ iterazioni.

Message Complexity: $O(n)$ per una iterazione, ripetuto per $n$ iterazioni
$$
\text{Messages}(\text{DistributedDijkstra}(G)) \in O(n^2)
$$

Time Complexity:
\begin{enumerate}
  \item init: catena di 2
  \item in ogni iterazione, solo l'invio di notifiche e ack è fatto in parallelo: \textit{start iteration}, risposta convergecast, notifica a $y_{\hat x}$, \textit{end iteration}.
  Totale di $\leq 4i$ messaggi in catena, ripetuti per $n-1$ iterazioni: $\leq 4 \sum_{i=1}^{n-1} i = 4 \frac{(n-1)n}{2} = 2(n-1)n$ messaggi.
  \item notifica di terminazione: catena di $\leq n-1$ messaggi
\end{enumerate}

Totale:
$$
\text{Time}(\text{DistributedDijkstra}(G)) = 2 + 2(n-1)n + (n-1) \in O(n^2)
$$

Distributed \underline{all} shortest path spanning trees: SPST costruito per ogni nodo. Costruire n SPST costa $O(n^3)$ messaggi.

Come si relazione con Gossip e Iterating (Distributed Bellman-Ford)? $O(mn) \in O(n^3)$ in Gossip e $O(mn^2) \in O(n^4)$ in Iterating.
Gossip gestisce meglio il ricalcolo delle tabelle di routing.
